
\documentclass[UTF8,a4paper,12pt,fontset=none]{ctexart}

\usepackage{geometry}
\usepackage{booktabs}
\usepackage{tabularx}
\usepackage{longtable}
\usepackage{amsmath}
\usepackage{amssymb}
\usepackage{graphicx}
\usepackage{xcolor}
\usepackage{hyperref}
\usepackage{enumitem}
\usepackage{float}
\usepackage{setspace}
\usepackage{caption}
\usepackage{titlesec}
\usepackage{tocloft}
\usepackage{array}
\usepackage{fontspec}

\setmainfont{DejaVu Sans}
\setsansfont{DejaVu Sans}
\setmonofont{Noto Sans Mono}
\setCJKmainfont{Noto Sans CJK SC}
\setCJKsansfont{Noto Sans CJK SC}
\setCJKmonofont{Noto Sans Mono CJK SC}

\geometry{left=3.05cm,right=3.05cm,top=3.0cm,bottom=3.0cm}
\emergencystretch=6em
\onehalfspacing
\setlength{\parskip}{0.42em}
\setlength{\parindent}{2em}
\renewcommand{\arraystretch}{1.22}
\setlength{\tabcolsep}{7pt}

\definecolor{aitInk}{HTML}{1D1D1F}
\definecolor{aitText}{HTML}{24313A}
\definecolor{aitBlue}{HTML}{0071E3}
\definecolor{aitGray}{HTML}{6E6E73}
\definecolor{aitRule}{HTML}{D2D2D7}
\definecolor{aitSoft}{HTML}{F5F5F7}
\definecolor{aitSoftBlue}{HTML}{EAF3FF}

\hypersetup{
  colorlinks=true,
  linkcolor=aitInk,
  urlcolor=aitBlue,
  citecolor=aitInk,
  pdfborder={0 0 0}
}

\setlist[itemize]{leftmargin=2em,itemsep=0.22em,topsep=0.40em}
\setlist[enumerate]{leftmargin=2em,itemsep=0.22em,topsep=0.40em}
\captionsetup{font=small,labelfont=bf,labelsep=quad}
\titleformat{\section}{\Large\bfseries\sffamily\color{aitInk}}{\thesection}{0.75em}{}
\titleformat{\subsection}{\large\bfseries\sffamily\color{aitText}}{\thesubsection}{0.75em}{}
\titlespacing*{\section}{0pt}{2.0em}{0.85em}
\titlespacing*{\subsection}{0pt}{1.2em}{0.5em}
\renewcommand{\cfttoctitlefont}{\Large\bfseries\sffamily\color{aitInk}}
\renewcommand{\cftsecfont}{\color{aitInk}}
\renewcommand{\cftsubsecfont}{\color{aitGray}}
\renewcommand{\cftsecpagefont}{\color{aitInk}}
\renewcommand{\cftsubsecpagefont}{\color{aitGray}}
\renewcommand{\cftsecleader}{\cftdotfill{\cftdotsep}}
\setlength{\cftbeforesecskip}{0.55em}
\setlength{\cftbeforesubsecskip}{0.12em}
\renewcommand{\cftdotsep}{1.6}

\newcolumntype{Y}{>{\raggedright\arraybackslash}X}
\newcolumntype{L}[1]{>{\raggedright\arraybackslash}p{#1}}

\newcommand{\metriccard}[3]{%
\begin{minipage}[t]{0.30\textwidth}
\setlength{\parindent}{0pt}
\colorbox{aitSoft}{\parbox{0.94\linewidth}{
\vspace{0.45em}
{\sffamily\bfseries\color{aitInk}#1}\par
\vspace{0.25em}
{\sffamily\fontsize{18}{22}\selectfont\color{aitBlue}#2}\par
\vspace{0.25em}
{\small\color{aitGray}#3}
\vspace{0.45em}
}}
\end{minipage}
}

\begin{document}

\begin{titlepage}
  \thispagestyle{empty}
  \vspace*{1.3cm}
  {\sffamily\small\color{aitGray}AItrader\hspace{0.8em}/\hspace{0.8em}MSGCA\par}
  \vspace{2.1cm}
  \noindent
  \begin{minipage}{0.80\textwidth}
    \raggedright
    {\sffamily\bfseries\fontsize{31}{39}\selectfont 基于多源因子与\\深度排序模型的\\A股Top20选股\\实验报告\par}
    \vspace{0.72cm}
    {\sffamily\large\color{aitGray}多源因子建模 · 横截面排序 · 主题上下文 · 动态聚类\par}
  \end{minipage}

  \vspace{1.75cm}
  \noindent
  \begin{minipage}{0.62\textwidth}
    {\color{aitRule}\rule{0.62\textwidth}{0.5pt}\par}
    \vspace{0.32cm}
    {\color{aitBlue}\rule{0.16\textwidth}{1.2pt}}
    \hspace{0.045\textwidth}
    {\color{aitGray}\rule{0.09\textwidth}{1.2pt}}
    \hspace{0.045\textwidth}
    {\color{aitInk}\rule{0.25\textwidth}{1.2pt}}\par
    \vspace{0.42cm}
    \hspace{0.02\textwidth}{\color{aitBlue}\rule{0.06\textwidth}{2.35cm}}
    \hspace{0.05\textwidth}{\color{aitSoftBlue}\rule{0.06\textwidth}{1.50cm}}
    \hspace{0.05\textwidth}{\color{aitInk}\rule{0.06\textwidth}{2.85cm}}
    \hspace{0.05\textwidth}{\color{aitRule}\rule{0.06\textwidth}{1.95cm}}
    \hspace{0.05\textwidth}{\color{aitBlue}\rule{0.06\textwidth}{2.60cm}}
  \end{minipage}

  \vfill
  \noindent{\color{aitRule}\rule{\textwidth}{0.5pt}\par}
  \vspace{0.45cm}
  \noindent
  \begin{tabularx}{\textwidth}{@{}L{2.8cm}Y@{}}
    {\color{aitGray}\small 课程项目} & {\small 股票趋势预测与模拟交易} \\
    {\color{aitGray}\small 小组成员} & {\small 待填写} \\
  \end{tabularx}
  \vspace*{0.8cm}
\end{titlepage}

\begin{center}
\begin{minipage}{0.88\textwidth}
\setlength{\parindent}{0pt}
\raggedright
\section*{摘要}

本实验面向 A 股日频数据构建短期 Top20 选股系统，目标是在严格时间可见性约束下，对每日股票横截面进行收益排序，并将排序结果转化为满仓轮换的模拟交易信号。系统以量价、基本面、资金流和新闻信息为输入，经过标准化、因子构造、样本对齐、因子筛选和深度学习建模后，输出每日个股 score。核心模型采用多源门控交叉注意力结构 MSGCA，将价格窗口、新闻因子与基本面/资金流/技术因子编码为不同模态 token，通过跨模态注意力和门控融合生成排序分数、收益预测和方向概率。

实验设计围绕“输入是否有效、训练目标是否匹配、模型结构是否稳定、比赛短窗口是否可解释”四个问题展开。早期消融表明，单一价格窗口的泛化不稳定，经过筛选的多源因子对横截面排序贡献更显著；后续结构实验发现，RankIC 提升并不必然带来 Top20 组合收益提升，过大的模型容量反而可能放大短窗口收益波动。进一步的 loss 和 score 实验说明，仅靠 TopK loss、分阶段训练或推理重排难以稳定解决收益目标错配问题。最终，引入主题上下文和动态聚类上下文后，模型在统一 holdout 窗口中获得更好的 period return 和 excess vs equal-weight。统一口径下，最终方案取得 12.56\% period return、16.36\% excess vs equal-weight、4.32\% rolling10 均值和 0.05848 competition score。

\end{minipage}
\end{center}

\vspace{0.6cm}
\tableofcontents
\newpage
\raggedright

\section{研究目标与任务定义}

课程项目要求在 A 股历史日频数据上构建包含深度学习模型的股票趋势预测系统，并基于预测结果完成历史回测与模拟交易。本实验将该任务抽象为每日开盘前的横截面排序问题：对目标交易日 $t$ 的股票 $i$，模型只能使用 $t$ 日之前已经可见的信息，输出排序分数 $s_{i,t}$。分数越高，表示模型认为该股票在下一交易窗口内的相对收益越高。

主监督标签定义为下一交易日开盘收益：
\[
  y_{i,t}=\frac{\mathrm{open}_{i,t+1}}{\mathrm{open}_{i,t}}-1.
\]
系统中该标签对应 \texttt{label\_next\_open\_return}。辅助标签为 \texttt{label\_next\_vwap\_return} 和收益方向标签。主标签贴近 T+1 交易制度下“今日形成信号、次日才能完整评估开盘收益”的比赛口径；VWAP 标签用于刻画实际执行价格与开盘价格之间的偏差。

实验目标不是单纯提高单点预测误差，而是优化“每日 Top20 股票组合”的排序质量。因此，本报告同时考察 RankIC、Top20 组合收益、excess vs equal-weight、rolling10 收益、latest10 收益和 max drawdown。RankIC 用于衡量横截面排序相关性，组合收益和超额收益用于评估交易端效果，rolling10 和 latest10 用于贴近约 10 个交易日的比赛周期。

\section{实验整体设计}

\subsection{实验主线}

本实验不是一组彼此孤立的模型试跑，而是按照问题诊断逐步推进。整体逻辑如表 \ref{tab:logic_chain} 所示。

\begin{table}[H]
  \centering
  \caption{实验推进逻辑}
  \label{tab:logic_chain}
  \begin{tabularx}{\textwidth}{L{0.18\textwidth}L{0.26\textwidth}Y}
    \toprule
    阶段 & 核心问题 & 设计动机 \\
    \midrule
    输入有效性检验 & 哪些信息源真正提供排序信号 & 先比较 price-only、factor-only、news+price 和 all-selected，判断模型应依赖价格窗口还是多源因子。 \\
    目标函数检验 & TopK loss 是否能直接提升组合收益 & 若模型已具备排序能力，则进一步检查训练目标是否与 Top20 组合收益一致。 \\
    结构稳定性检验 & prototype、gate 和模型容量如何影响交易表现 & 当 RankIC 与收益出现不一致时，需要判断问题来自结构表达、门控机制还是容量过大。 \\
    策略输出检验 & 模型分数如何转化为可交易组合 & 比较 TopN、每日换仓数、收益头、分数融合、市值过滤和推理重排，判断模型输出到交易信号之间的稳定转换方式。 \\
    短窗口对齐检验 & 模型是否贴近 10 日比赛周期 & 使用 rolling10 监督、多头分阶段训练和 score 变体，检验训练目标与部署指标的匹配程度。 \\
    市场状态建模 & 主题与细分股票簇是否解释近期收益 & 在常规模型无法解释近期窗口差异时，引入主题上下文和动态聚类上下文。 \\
    \bottomrule
  \end{tabularx}
\end{table}

上述设计使报告形成从基础输入到最终部署方案的递进关系：先确认多源因子确实优于单一价格窗口，再确认简单调 loss 或扩大模型并不能稳定改善收益，最后把改进重点转向市场状态和细分主题结构。

\subsection{核心指标}

\begin{center}
\metriccard{最终 period return}{12.56\%}{统一 holdout 窗口中的 Top20 组合收益}
\hfill
\metriccard{excess vs equal-weight}{16.36\%}{组合收益减去同期等权市场收益}
\hfill
\metriccard{rolling10 均值}{4.32\%}{贴近比赛周期的 10 日滚动收益}
\end{center}

\vspace{0.4cm}

\begin{table}[H]
  \centering
  \caption{评估指标说明}
  \label{tab:metric_guide}
  \begin{tabularx}{\textwidth}{L{0.23\textwidth}Y}
    \toprule
    指标 & 含义 \\
    \midrule
    RankIC & 每个交易日模型分数与真实收益的横截面 Spearman 相关性，再按交易日平均。该指标衡量整体排序方向是否正确。 \\
    组合收益 & 按模型分数构造 Top20 等权组合后的 period return，反映策略端绝对表现。 \\
    excess vs equal-weight & Top20 组合收益减去同期全市场等权组合收益，用于区分模型收益与市场整体 beta。 \\
    rolling10 & 每 10 个交易日滚动窗口的 Top20 组合收益均值，用于贴近比赛周期。 \\
    latest10 & 最近 10 个交易日窗口收益，用于观察模型在最新市场状态下的表现。 \\
    max drawdown & 组合净值从阶段高点到阶段低点的最大跌幅，用于衡量风险暴露。 \\
    competition score & 实验内部用于比较方案的综合指标，融合 period return、rolling10、回撤等维度。 \\
    \bottomrule
  \end{tabularx}
\end{table}

\section{数据处理与防泄露约束}

\subsection{数据来源与样本构造}

实验使用课程提供的 A 股日频数据，不额外引入外部市场数据。数据类型包括日频量价、基本面指标、资金流、新闻文本、ST 状态和股票基础信息。数据处理流程分为四层：

\begin{enumerate}
  \item 原始表清洗：统一股票代码、交易日期、字段类型和缺失值处理方式。
  \item 标准宽表构建：将价格、基本面、资金流和 ST 状态按股票和日期对齐。
  \item 因子构造：生成价格形态、动量反转、波动率、估值、市值、资金流、新闻情绪和市场状态类因子。
  \item 样本级对齐：把日频因子按决策时点对齐到样本，形成模型输入和监督标签。
\end{enumerate}

样本以“股票—目标交易日”为基本单位。每条样本记录目标交易日、特征可见日期、主标签、辅助标签和特征 block。特征 block 覆盖价格窗口、新闻聚合、技术因子、资金流因子、基本面因子和 GPT 挖掘因子。实验中使用两类特征清单：严格消融阶段使用较早筛选版本，主要用于时间切分下的结构比较；比赛复训阶段使用更新后的 reviewed 特征清单，主要用于贴近最新市场状态的部署测试。

\subsection{时间切分}

为了避免未来信息进入训练过程，所有实验采用时间顺序切分，而不是随机切分。默认切分如下：

\begin{table}[H]
  \centering
  \caption{数据切分设置}
  \begin{tabularx}{\textwidth}{L{0.22\textwidth}L{0.28\textwidth}Y}
    \toprule
    数据集 & 时间范围 & 用途 \\
    \midrule
    train set & 2019年1月1日至2025年9月30日 & 拟合模型参数和标准化统计量。 \\
    validation set & 2025年10月1日至2025年12月31日 & 模型选择、超参数比较和早期结构诊断。 \\
    holdout & 2026年1月1日之后 & 检验模型在未参与训练区间的排序和组合表现。 \\
    统一 holdout 窗口 & 2026年5月6日至2026年6月1日 & 对后期框架变体进行统一比较，贴近比赛前市场环境。 \\
    \bottomrule
  \end{tabularx}
\end{table}

其中，比赛复训模型会使用更靠近比赛期的数据，但这类实验仅作为部署方案选择，不作为严格独立测试结论。

\subsection{防泄露策略}

金融时间序列实验最主要的风险是数据泄露。为保证实验结论可解释，本实验设置以下约束：

\begin{itemize}
  \item 每条样本只使用不晚于 \texttt{feature\_asof\_date} 的行情、资金流、基本面和新闻信息。
  \item 滚动窗口类特征只沿历史方向计算，不使用目标日及之后的数据。
  \item 标准化器仅在 train set 拟合，validation、holdout 和推理阶段只复用 train set 统计量。
  \item 训练和评估均按交易日组织 batch，避免同一日期的横截面标签被错误打乱。
  \item 新闻评分先离线转化为结构化因子，再进入训练流程；训练阶段不动态调用外部语言模型。
  \item ST 股票和不符合比赛范围的股票在回测与推理阶段采用一致过滤口径。
\end{itemize}

早期实验中曾出现收益异常偏高的结果。核查后发现该版本包含目标日新闻统计字段，存在未来信息风险。因此该结果不进入正式对比，仅作为防泄露检查的反例。这一处理保证了后续实验比较建立在一致、可审计的输入边界之上。

\section{特征工程}

\subsection{人工因子}

人工因子覆盖四类信息。第一类为价格与成交类因子，包括收益率、均线、成交量变化、波动率、振幅、动量和反转。第二类为基本面与估值因子，包括市值、PE、PB、换手率和股息率等字段及其横截面派生。第三类为资金流因子，包括主力、小单、大单和超大单资金流强度及其滚动变化。第四类为行业中性化因子，通过同一交易日和同一行业内的均值扣除，减少行业整体涨跌对个股排序的干扰。

因子构造同时包含时序算子和横截面算子。时序算子按股票分组并沿交易日方向计算，用于捕捉个股自身历史状态；横截面算子按交易日分组，用于捕捉同日股票之间的相对位置。该设计符合 Top20 选股任务的本质：模型最终比较的是同一交易日中不同股票的相对吸引力，而不是单只股票的绝对价格水平。

\subsection{新闻与文本因子}

新闻数据先通过离线语言模型评分，转化为情绪、关注度和潜在影响类因子，再按股票和日期聚合到样本级。新闻因子不直接作为原始文本输入模型，而是作为结构化特征参与排序。这样做有两个原因：一是降低训练成本，避免每次训练重复处理长文本；二是保持时间可见性清晰，避免在训练过程中引入不可控的在线调用结果。

实验中也使用 GPT 挖掘备选因子，主要覆盖行业相对强弱、流动性、资金流交互、市场状态、反转、风险控制和估值等方向。这些因子必须通过字段合法性检查、时间可见性检查和缺失率检查后才能进入候选特征池。最终 reviewed 特征清单保留了具有一定经济含义和统计质量的因子，同时剔除明显高缺失、高相关冲突或解释性较弱的字段。

\section{模型方法}

\subsection{MSGCA 结构}

本实验主模型为 MSGCA，即多源门控交叉注意力模型。模型输入包含三路信息：

\begin{itemize}
  \item 价格窗口：最近 60 个交易日的 OHLC、VWAP、成交量、成交额和资金流比例等变量。
  \item 新闻与文本因子：新闻数量、市场新闻状态、情绪评分和影响强度聚合特征。
  \item 基本面/资金流/技术因子：经过筛选的 Alpha、估值、市值、资金流和 GPT 挖掘因子。
\end{itemize}

价格窗口采用变量 token 表示。每个价格变量被视为一个时间序列 token，经归一化和 Transformer 编码后得到价格表征。数值因子采用 factor-aware encoder：每个标量因子先映射为 feature token，再结合 feature embedding 和 group embedding；组内通过 factor gate 聚合为 group token，从而保留因子分组结构和组内权重信息。

\subsection{跨模态融合}

设价格、新闻和因子三路 token 分别为 $H_p,H_t,H_f$。模型通过跨模态注意力让价格 token 关注新闻与因子上下文，同时让因子 token 感知近期价格状态。随后对各模态 token 池化，得到 $z_p,z_t,z_f$。门控网络输出三路模态权重：
\[
  (g_p,g_t,g_f)=\operatorname{softmax}(W[z_p,z_t,z_f,m_p,m_t,m_f]),
\]
其中 $m_p,m_t,m_f$ 表示模态可用性。最终融合向量为：
\[
  z=g_pz_p+g_tz_t+g_fz_f.
\]

模型输出三个头：
\[
  s=\operatorname{Head}_{rank}(z),\quad
  \hat{r}=\operatorname{Head}_{return}(z),\quad
  \hat{d}=\operatorname{Head}_{direction}(z).
\]
其中 $s$ 用于每日横截面排序，$\hat{r}$ 用于收益回归辅助学习，$\hat{d}$ 用于上涨方向分类辅助学习。

\subsection{训练目标}

模型训练损失由排序损失、收益回归损失、方向分类损失、TopK 收益近似损失和门控熵正则组成：
\[
\mathcal{L}
=
\lambda_{rank}\mathcal{L}_{rank}
+\lambda_{ret}\mathcal{L}_{mse}
+\lambda_{dir}\mathcal{L}_{bce}
+\lambda_{topk}\mathcal{L}_{topk}
-\lambda_{ent}\mathcal{H}(g).
\]

其中，$\mathcal{L}_{rank}$ 为按交易日分组的 LambdaRank 风格排序损失，优化同一交易日横截面内的相对顺序；$\mathcal{L}_{topk}$ 使用 softmax 权重近似 TopK 组合收益目标；门控熵正则用于避免模型过早塌缩到单一模态。该目标同时考虑排序、收益方向和交易端 TopK 选择，但实验结果显示，多目标训练仍需要与具体市场状态和推理分数设计配合，不能仅依赖单个 loss 权重调节。

\section{回测与交易规则}

回测采用与比赛接近的 Top20 满仓轮换策略。初始资金为 100 万元，第一日按模型分数等权买入前 20 只股票。之后每日从可卖持仓中卖出分数较低的 3 只，同时从全市场买入分数较高且当前未持有的 3 只，尽量保持满仓。回测包含手续费、滑点和 T+1 卖出限制。

该策略设计有三个目的。第一，Top20 等权持仓与比赛“选股组合”目标一致。第二，每日只轮换部分股票，避免模型分数微小变化导致过度交易。第三，加入 T+1 限制、手续费和滑点后，回测收益更接近真实模拟交易环境。需要注意的是，回测仍然是简化模型，尚未完整刻画涨跌停、停牌、盘口深度和委托成交概率等约束，因此回测结果应被解释为模型信号质量的近似评估，而不是实盘收益保证。

\section{策略输出实验}

模型训练得到的是每日个股分数，而比赛提交需要的是可执行的买入列表。策略输出实验只改变分数到组合的映射方式，不改变模型权重。为避免把网格搜索结果直接堆叠到正文，本节按实验目的分为四类：回测口径、组合构造、分数融合、约束与重排。

\begin{table}[H]
  \centering
  \caption{策略输出实验分类}
  \begin{tabularx}{\textwidth}{L{0.20\textwidth}L{0.30\textwidth}Y}
    \toprule
    类别 & 主要变量 & 统计目的 \\
    \midrule
    回测口径校验 & holdout 窗口、10 日窗口、最新短窗口 & 判断回测流程是否稳定，并确认最终评估不能只依赖长区间平均收益。 \\
    组合构造网格 & TopN、每日换仓数、T+1 换手约束 & 比较集中持仓与分散持仓的收益和风险差异。 \\
    分数融合网格 & \texttt{y\_score}、\texttt{return\_pred}、blend、direction & 判断排序头、收益头和方向头在交易端分别提供什么信息。 \\
    约束与重排网格 & 市值过滤、cap bonus、cluster boost、固定簇筛选 & 检查推理阶段规则是否能改善可交易性和短窗口稳定性。 \\
    \bottomrule
  \end{tabularx}
\end{table}

\subsection{分类统计摘要}

\begin{table}[H]
  \centering
  \caption{策略输出实验代表性结果}
  \resizebox{\textwidth}{!}{
  \begin{tabular}{llrrrrr}
    \toprule
    类别 & 代表变体 & TopN / 换仓 & period return & excess vs equal-weight & rolling10 或 latest & max drawdown \\
    \midrule
    回测口径 & factor-aware holdout & 20 / 3 & 7.30\% & -- & -- & -13.27\% \\
    回测口径 & 2026-05-21 短窗口 & 20 / 3 & -4.77\% & -- & -- & -5.77\% \\
    组合构造 & \texttt{z\_y\_plus\_return} & 5 / 0 & 8.93\% & 13.12\% & -- & -- \\
    组合构造 & \texttt{y\_score} & 20 / 3 & -5.16\% & -0.97\% & 1.21\% & -6.15\% \\
    分数融合 & \texttt{return\_pred} & 5 / 1 & 50.30\% & 46.15\% & latest -4.66\% & -20.03\% \\
    分数融合 & \texttt{blend\_y25\_ret75} & 5 / 1 & 46.67\% & 42.53\% & latest -2.03\% & -18.17\% \\
    约束规则 & \texttt{return\_pred + cap} & 20 / 5 & 20.94\% & 9.04\% & latest 1.31\% & -8.70\% \\
    推理重排 & \texttt{cluster\_boostlight} & 20 / 规则 & 5.33\% & 9.14\% & rolling 1.64\% & -6.21\% \\
    推理重排 & 固定 cluster 筛选 & 20 / 规则 & -7.82\% 至 -13.86\% & -4.01\% 至 -10.06\% & -- & -15.32\% 至 -18.51\% \\
    \bottomrule
  \end{tabular}
  }
\end{table}

观察与分析：
\begin{itemize}
  \item 回测口径校验显示，完整 holdout 和最后几天短窗口的收益方向可以相反。因此后续评估加入 rolling10、latest10 和 excess vs equal-weight，避免仅按单个长区间选模型。
  \item 组合构造网格显示，Top5/Top10 集中组合在部分窗口收益更高，但 max drawdown 和窗口波动也更高；Top20 更符合比赛满仓组合目标，也更适合作为最终部署口径。
  \item 分数融合网格显示，\texttt{return\_pred} 在交易层具有明显信息量，不能只把它当作训练辅助头。但 \texttt{return\_pred} 主导的 Top5 策略回撤较大，说明它适合作为训练和 score 设计依据，不适合作为无约束集中持仓规则。
  \item 市值 cap 和轻量 cluster boost 能改善部分窗口表现，但固定簇数硬筛选明显降低收益。该结果支持后续采用“聚类信息进入训练上下文”，而不是在推理阶段强行按簇筛股票。
\end{itemize}

上述策略实验的作用是把问题从“模型分数是否有效”进一步定位到“哪一种模型输出最适合服务 Top20 组合”。因此，后续模型实验不只比较 RankIC，还把 rolling10、latest10、excess vs equal-weight 和推理层 score 一起纳入评估。

\section{实验结果与分析}

\subsection{实验一：输入与特征消融}

第一组实验检验不同信息源对排序能力和组合收益的贡献。模型结构基本保持一致，主要改变输入模态和特征集合。

\begin{table}[H]
  \centering
  \caption{输入与特征消融结果}
  \label{tab:ablation}
  \resizebox{\textwidth}{!}{
  \begin{tabular}{lrrrrrrr}
    \toprule
    局部变体 & validation RankIC & validation return & validation Sharpe & holdout RankIC & holdout return & holdout excess & holdout Sharpe \\
    \midrule
    price\_only & 0.0406 & 7.29\% & 1.75 & 0.0208 & -3.61\% & -15.08\% & -0.39 \\
    factor\_only & 0.0734 & 4.04\% & 1.10 & 0.0450 & 6.92\% & -4.55\% & 0.97 \\
    manual\_price & 0.0665 & 3.93\% & 1.09 & 0.0415 & 4.17\% & -7.30\% & 0.64 \\
    news\_price & 0.0444 & 1.27\% & 0.39 & 0.0227 & -0.30\% & -11.77\% & 0.07 \\
    all\_selected & 0.0614 & 6.04\% & 1.63 & 0.0414 & 7.30\% & -4.16\% & 1.03 \\
    factor\_aware\_h48 & 0.0680 & 5.10\% & 1.27 & 0.0424 & 11.86\% & 0.40\% & 1.59 \\
    \bottomrule
  \end{tabular}
  }
\end{table}

结果显示，\texttt{price\_only} 在 validation 上表现尚可，但 holdout 收益转负，说明单一价格窗口容易受到市场风格变化影响。相比之下，\texttt{factor\_only} 的 holdout RankIC 明显更高，表明经过筛选的多源因子对横截面排序有稳定贡献。加入 factor-aware 编码后，\texttt{factor\_aware\_h48} 在 holdout 取得 11.86\% 组合收益和 0.40\% excess vs equal-weight，是该组中交易端表现最好的方案。

该实验给出后续设计的基础结论：模型不能只依赖价格窗口，应把多源因子作为主要输入；同时，评价模型不能只看 validation 收益，还必须检查 holdout 表现和 excess vs equal-weight。

\subsection{实验二：TopK Loss 权重消融}

第二组实验检验 TopK loss 权重是否能直接提升 Top20 组合收益。模型结构固定为 factor-aware h64，仅改变 $\mathcal{L}_{topk}$ 权重。

\begin{table}[H]
  \centering
  \caption{TopK loss 权重消融}
  \label{tab:topk_ablation}
  \begin{tabular}{lrrr}
    \toprule
    局部变体 & RankIC & 组合收益 & excess vs equal-weight \\
    \midrule
    topk000 & 0.0415 & 7.24\% & -4.23\% \\
    topk002 & 0.0414 & 7.30\% & -4.16\% \\
    topk005 & 0.0413 & 6.41\% & -5.05\% \\
    \bottomrule
  \end{tabular}
\end{table}

\texttt{topk000} 和 \texttt{topk002} 的结果接近，说明小权重 TopK loss 没有显著改变整体排序结构；当权重提高到 0.05 时，组合收益和超额收益反而下降。该结果说明，TopK loss 并不是越大越好。过强的 TopK 约束可能使模型过度关注训练期局部收益最高的样本，从而削弱 holdout 稳定性。因此，后续实验不再简单提高 TopK loss，而是转向模型内部结构和容量诊断。

\subsection{实验三：Prototype 与 Gate 消融}

第三组实验检验因子分组 prototype 和 factor gate 对收益表现的影响。该组实验的核心问题是：当 RankIC 与组合收益不一致时，结构模块是否在其中发挥作用。

\begin{table}[H]
  \centering
  \caption{Prototype 与 factor gate 消融}
  \label{tab:prototype_gate_ablation}
  \begin{tabular}{lrrr}
    \toprule
    局部变体 & RankIC & 组合收益 & excess vs equal-weight \\
    \midrule
    plain & 0.0413 & 6.41\% & -5.05\% \\
    proto2 & 0.0427 & 5.79\% & -5.68\% \\
    no\_gate & 0.0454 & -0.75\% & -12.22\% \\
    \bottomrule
  \end{tabular}
\end{table}

加入 prototype 后，RankIC 有小幅提升，但组合收益没有同步改善；移除 gate 后，RankIC 最高，组合收益却转为负值。这说明横截面整体相关性的提高不等价于 Top20 组合收益提高。尤其是 \texttt{no\_gate} 的结果表明，gate 虽然可能限制部分排序相关性，但对交易端风险有保护作用。没有 gate 时，模型更容易放大噪声因子，使高分股票集中到短窗口中不稳定的样本。

该实验把问题进一步明确为：模型不是不会排序，而是整体排序指标与 Top20 组合目标存在错配。后续实验需要关注模型容量和短窗口目标，而不是仅追求更高 RankIC。

\subsection{实验四：Hidden Size 消融}

第四组实验检验模型容量对收益稳定性的影响。结构固定为 proto2/topk005，仅改变隐藏维度。

\begin{table}[H]
  \centering
  \caption{Hidden size 消融}
  \label{tab:hidden_size_ablation}
  \begin{tabular}{lrrr}
    \toprule
    局部变体 & RankIC & 组合收益 & excess vs equal-weight \\
    \midrule
    h48 & 0.0424 & 11.86\% & 0.40\% \\
    h64 & 0.0427 & 5.79\% & -5.68\% \\
    h96 & 0.0432 & -0.37\% & -11.83\% \\
    \bottomrule
  \end{tabular}
\end{table}

随着 hidden size 从 48 增大到 96，RankIC 略有上升，但组合收益持续恶化。这说明更大的模型确实学习到了某些对整体横截面排序有帮助的细节，但这些细节没有转化为 Top20 组合收益。对于短期 A 股选股任务而言，收益标签噪声较高，模型容量过大时更容易捕捉训练期局部模式，导致 holdout 中 Top20 高分股票质量下降。

因此，后续实验选择 h48 作为主要容量设置。这一选择不是因为 h48 的表达能力最强，而是因为它在当前数据规模、标签噪声和交易目标下具有更好的收益稳定性。

\subsection{实验五：最新特征复训与 Gate 函数对比}

在确定小模型更稳后，第五组实验使用更新后的 reviewed 特征清单进行复训，并比较 sparsemax 与 softmax gate。该实验用于判断近期特征更新和门控函数形式是否能够改善短窗口收益。

\begin{table}[H]
  \centering
  \caption{最新特征复训与 Gate 函数对比}
  \resizebox{\textwidth}{!}{
  \begin{tabular}{llrrrr}
    \toprule
    局部变体 & Gate & RankIC & 组合收益 & excess vs equal-weight & Sharpe \\
    \midrule
    sparse & sparsemax & -0.0472 & -1.81\% & -4.04\% & -2.29 \\
    soft & softmax & -0.0479 & -1.69\% & -3.92\% & -2.55 \\
    soft\_recent & softmax & -0.0479 & -1.69\% & -3.92\% & -2.55 \\
    \bottomrule
  \end{tabular}
  }
\end{table}

三种设置的 RankIC 均为负，组合收益也为负。softmax 相比 sparsemax 的亏损略小，但差异不足以解释整体表现下降。加入近期样本权重后，结果没有明显变化。这表明短窗口收益下降不是单纯由 gate 函数形式或最新特征数量决定的，更可能来自训练目标与当前市场状态之间的错配。

该结论促使后续实验将重点从“换特征、换 gate”转向“让训练和推理更贴近比赛窗口”。

\subsection{实验六：多头分阶段训练与 Rolling-10 监督}

第六组实验针对多头输出不一致问题进行改造。原始 MSGCA 同时输出排序分数、收益预测和方向概率，但实际交易只使用最终排序分数。分阶段训练的思路是先分别校准多头输出，再联合优化 final score，并加入 rolling10 监督，使训练目标更接近 10 个交易日的比赛周期。

\begin{table}[H]
  \centering
  \caption{多头分阶段训练结果}
  \resizebox{\textwidth}{!}{
  \begin{tabular}{lrrr}
    \toprule
    局部变体 & period return & excess vs equal-weight & max drawdown \\
    \midrule
    stagewise\_final & -7.26\% & -0.35\% & -8.20\% \\
    swloss\_v3 & -0.42\% & 3.38\% & -8.28\% \\
    \bottomrule
  \end{tabular}
  }
\end{table}

\texttt{stagewise\_final} 的 period return 为负，说明分阶段校准多头输出并不能直接提升短窗口 Top20 收益。\texttt{swloss\_v3} 的 excess vs equal-weight 为正，但绝对收益仍然较弱，说明 rolling10 监督可以在一定程度上减少市场 beta 干扰，却仍不足以产生稳定的绝对收益。

该实验说明，改进 loss 可以改善部分相对表现，但不能单独解决市场状态变化问题。模型需要能够识别当期更有效的主题和股票簇结构。

\subsection{实验七：严格字段约束下的 Loss/Score 方案}

第七组实验将 loss 和推理 score 分开测试，并严格限制在已经物化的字段内实现，不引入尚未落地的行业、主题或资金流派生变量。

\begin{table}[H]
  \centering
  \caption{严格字段约束下的 loss/score 方案}
  \resizebox{\textwidth}{!}{
  \begin{tabular}{lrrrrrr}
    \toprule
    局部变体 & period return & excess vs equal-weight & rolling10 & latest10 & max drawdown & competition score \\
    \midrule
    l1\_multihead & 2.57\% & 6.38\% & 4.35\% & 0.08\% & -7.02\% & 0.04838 \\
    exact\_news\_final & 5.71\% & 9.52\% & 3.94\% & -3.05\% & -7.97\% & 0.04528 \\
    exact\_news\_score & 6.51\% & 10.32\% & 2.99\% & 1.41\% & -8.68\% & 0.03876 \\
    s5\_rerank & 1.48\% & 5.29\% & -1.07\% & -0.45\% & -8.27\% & -0.00433 \\
    \bottomrule
  \end{tabular}
  }
\end{table}

\texttt{l1\_multihead} 的 competition score 最高，主要来自较好的 rolling10 表现；\texttt{exact\_news\_score} 的 period return 更高，但 rolling10 和 max drawdown 更弱，说明单看整段收益会高估该方案。复杂重排方案 \texttt{s5\_rerank} 表现较差，说明在现有字段约束下，推理阶段手工重排没有稳定增益。

这一组实验的作用是排除“仅通过已有字段组合就能解决问题”的可能性。由于严格字段方案没有充分解释近期主题行情，下一步需要显式建模主题上下文。

\subsection{实验八：主题上下文训练}

第八组实验在模型输入中加入主题上下文 token，包括主题动量、资金确认和回撤健康度等训练期可见特征。该设计的目的不是在推理阶段手工指定行业，而是让模型在训练中学习主题强度对个股收益的影响。

\begin{table}[H]
  \centering
  \caption{主题上下文训练结果}
  \resizebox{\textwidth}{!}{
  \begin{tabular}{lrrrrrr}
    \toprule
    局部变体 & period return & excess vs equal-weight & rolling10 & latest10 & max drawdown & competition score \\
    \midrule
    ctx002\_medium & 10.89\% & 14.70\% & 5.00\% & 4.33\% & -5.59\% & 0.06332 \\
    ctx005\_medium & 2.42\% & 6.22\% & 2.88\% & -7.57\% & -11.95\% & 0.02854 \\
    ctx008\_medium & 11.54\% & 15.35\% & 5.02\% & 4.62\% & -4.64\% & 0.06506 \\
    ctx008\_soft & 12.27\% & 16.08\% & 2.91\% & 2.39\% & -4.97\% & 0.04760 \\
    \bottomrule
  \end{tabular}
  }
\end{table}

主题上下文显著改善了统一 holdout 窗口表现。其中，\texttt{ctx008\_medium} 的 competition score 最高，rolling10 和 latest10 也较稳定；\texttt{ctx008\_soft} 的 period return 最高，但 rolling10 相对较弱。该结果说明，主题强度对近期 Top20 选股有明显解释力，但 score 组合需要在收益峰值和窗口稳定性之间权衡。

不过，主题上下文仍然依赖预设主题字段，难以充分表达粗行业内部的细分链条。为进一步刻画更细粒度的股票共振结构，下一步引入动态聚类上下文。

\subsection{实验九：动态聚类上下文训练}

第九组实验使用动态聚类刻画细分股票簇。聚类信息不作为推理阶段硬筛规则，而是在训练样本中提供簇编号、簇强度、簇内资金确认和簇内排序约束，使模型在共享表示中学习细分主题结构。

\begin{table}[H]
  \centering
  \caption{主要框架变体的统一 holdout 指标}
  \label{tab:late_stage_metrics}
  \resizebox{\textwidth}{!}{
  \begin{tabular}{lrrrrrr}
    \toprule
    方案 & period return & excess vs equal-weight & rolling10 & latest10 & max drawdown & competition score \\
    \midrule
    stagewise\_final & -7.26\% & -0.35\% & -- & -- & -8.20\% & -- \\
    l1\_multihead & 2.57\% & 6.38\% & 4.35\% & 0.08\% & -7.02\% & 0.04838 \\
    theme\_ctx008\_medium & 11.54\% & 15.35\% & 5.02\% & 4.62\% & -4.64\% & 0.06506 \\
    baseline\_medium & 7.73\% & 11.53\% & 3.19\% & 4.13\% & -5.13\% & 0.04517 \\
    cluster\_inrank020\_soft & 12.56\% & 16.36\% & 4.32\% & 0.70\% & -5.68\% & 0.05848 \\
    cluster\_ensemble\_soft & 7.34\% & 11.15\% & 3.43\% & -0.39\% & -6.75\% & 0.04403 \\
    \bottomrule
  \end{tabular}
  }
\end{table}

\texttt{cluster\_inrank020\_soft} 取得最高 period return 和最高 excess vs equal-weight，说明动态聚类上下文改善了统一 holdout 窗口中的 Top20 组合表现。与主题上下文相比，cluster 版本的 latest10 较弱，但整段收益和超额收益更高，体现出更强的收益捕捉能力。多 seed ensemble 反而低于单模型，说明该任务的 Top20 组合需要保留高分股票的集中度，简单平均会削弱模型的信号尖锐度。

推理层诊断进一步显示，hard cluster boost 和固定簇数筛选均低于直接使用 soft score。这说明聚类信息更适合作为训练上下文，而不是推理阶段硬约束。最终方案因此采用动态聚类上下文训练，但不叠加硬聚类筛选规则。

\section{综合讨论}

\subsection{从实验结果得到的主要结论}

第一，多源因子是本任务的核心输入。单一价格窗口在 validation 上可以获得一定收益，但 holdout 表现不稳定；经过筛选的基本面、资金流、技术和新闻因子能提供更稳定的横截面排序信息。

第二，RankIC 与 Top20 组合收益并不完全一致。若模型提高的是整体横截面排序相关性，但提升主要来自中位股票排序，则未必能改善 Top20 组合收益。实验中 \texttt{no\_gate} 和 h96 均体现了这一问题：RankIC 更高，但组合收益更差。

第三，模型容量需要受控。短期股票收益噪声较高，扩大 hidden size 会提高表达能力，但也更容易学习训练期局部结构。h48 的表现说明，在该任务中更小的模型反而更适合作为稳定部署基线。

第四，单纯调整 loss 或推理 score 不足以解决短窗口收益问题。TopK loss、分阶段训练、rolling10 监督和严格字段 score 都能改善部分指标，但无法稳定解释近期市场状态。更关键的改进来自主题上下文和动态聚类上下文。

第五，主题和细分股票簇对近期比赛窗口具有明显解释力。主题上下文提升了 rolling10 和 latest10 稳定性，动态聚类上下文提升了 period return 和 excess vs equal-weight。这说明模型需要捕捉市场资金在细分方向上的集中行为，而不是只依赖个股静态因子。

\subsection{最终方案选择}

综合 period return、excess vs equal-weight、rolling10、max drawdown 和窗口稳定性，本实验选择动态聚类上下文训练模型作为最终方案。该方案的优势在于：

\begin{itemize}
  \item 在统一 holdout 窗口中取得最高 period return 和最高 excess vs equal-weight。
  \item 聚类信息进入训练过程，而不是作为推理阶段硬筛规则，降低了过度人工干预的风险。
  \item 相比单纯 score 重排，动态聚类能在表示层刻画细分主题结构，更符合近期 A 股主题轮动特征。
  \item 最终推理仍保留连续分数排序，便于与 Top20 满仓轮换策略衔接。
\end{itemize}

需要同时注意，主题上下文方案在 competition score 和 latest10 上具有竞争力。因此最终部署时应持续监控动态聚类方案与主题上下文方案的最新窗口表现。如果市场从集中主题行情切换为更分散的风格，主题上下文或小模型基线可能重新具有优势。

\section{局限性与改进方向}

本实验仍存在以下限制。

第一，评估窗口较短。统一 holdout 窗口接近比赛环境，但其长度有限，不能等同于长期稳定收益证明。后续应在更长时间区间上做滚动 walk-forward 验证。

第二，回测成交假设仍然简化。当前回测包含手续费、滑点和 T+1 限制，但尚未完整模拟涨跌停、停牌、盘口深度、实际可成交量和委托价格偏离。若用于更严格的实盘评估，需要加入更细的交易撮合约束。

第三，新闻因子依赖离线评分质量。新闻评分可以压缩文本信息，但也可能丢失事件细节。后续可以比较结构化新闻因子、文本 embedding 和事件抽取特征的差异。

第四，动态聚类仍有过拟合风险。聚类上下文在当前 holdout 窗口有效，但主题簇结构可能随市场状态变化。后续应加入聚类稳定性评价，例如簇内收益一致性、簇切换频率和簇强度衰减监控。

第五，最终 score 仍以排序为主，尚未显式建模组合层风险。后续可以加入组合相关性、行业暴露、波动率约束和持仓换手约束，使模型从“选出最高分股票”进一步转向“选出风险调整后更优的 Top20 组合”。

\section{结论}

本实验完成了从数据清洗、因子构造、深度学习排序建模到模拟交易回测的完整流程。与直接训练价格序列模型相比，本系统的主要改进在于：构建了严格的时间可见性约束；将价格、新闻、基本面、资金流和技术因子统一对齐到样本级；使用 factor-aware MSGCA 建模多源信息；通过 Top20 满仓轮换策略评估模型在交易端的实际效果。

实验结果表明，多源因子优于单一价格窗口，小模型优于过大容量模型，单纯提高 RankIC 或增强 TopK loss 不一定改善组合收益。真正带来后期提升的是对市场状态的建模：主题上下文增强了模型对近期主题行情的识别能力，动态聚类上下文进一步刻画了粗行业内部的细分股票簇结构。最终方案在统一 holdout 窗口中取得 12.56\% period return、16.36\% excess vs equal-weight、4.32\% rolling10 均值和 0.05848 competition score。

总体而言，本实验的核心结论是：A 股短期 Top20 选股不能只依赖更复杂的网络结构或更激进的 loss，而需要在严格防泄露的基础上，将多源因子、交易目标和近期市场状态统一建模。动态聚类上下文训练为这一目标提供了当前实验中最有效的实现方式。

\appendix

\section{附录：主要实验矩阵}

\small
\begin{longtable}{L{0.15\textwidth}L{0.23\textwidth}L{0.18\textwidth}L{0.25\textwidth}}
  \caption{实验模块与结论汇总} \\
  \toprule
  实验模块 & 主要比较 & 最优或代表方案 & 结论 \\
  \midrule
  \endfirsthead
  \toprule
  实验模块 & 主要比较 & 最优或代表方案 & 结论 \\
  \midrule
  \endhead
  输入消融 & price-only、factor-only、news+price、all-selected & factor-aware h48 & 多源因子比单一价格窗口更稳定。 \\
  TopK loss & 权重 0、0.02、0.05 & topk002 接近 topk000 & 提高 TopK loss 不能直接提升收益。 \\
  结构消融 & plain、proto2、no-gate & 保留 gate & gate 对交易端有风险保护作用。 \\
  容量消融 & h48、h64、h96 & h48 & 更大模型带来更高过拟合风险。 \\
  最新复训 & sparsemax、softmax、近期样本权重 & 无明显最优 & gate 函数不是短窗口下降主因。 \\
  分阶段训练 & stagewise final、swloss v3 & swloss v3 & rolling10 监督改善相对收益，但绝对收益不足。 \\
  严格字段方案 & 多头 loss、新闻 score、重排 score & l1 multihead & 已有字段组合无法充分解释主题行情。 \\
  主题上下文 & 不同主题强度与 score 组合 & ctx008 medium & 主题状态能显著改善短窗口表现。 \\
  动态聚类 & 聚类上下文、ensemble、硬筛诊断 & cluster inrank020 soft & 聚类更适合作为训练上下文，而非推理硬规则。 \\
  \bottomrule
\end{longtable}
\normalsize

\end{document}
